\section{Experiments}

\subsection{Dataset and Protocol}
We use six VisA categories required by the project: candle, capsules, cashew,
chewinggum, fryum, and pipe\_fryum. Table~\ref{tab:dataset} summarizes the
official 1-class split. The original 1.9 GB dataset remains read-only. We build
the AnomalyDINO directory layout with symlinks and store only metadata,
selected references, prediction maps, and results.

\begin{table}[t]
\centering
\small
\caption{VisA categories and official 1-class test counts.}
\label{tab:dataset}
\begin{tabular}{lrrr}
\toprule
Category & Train normal & Test normal & Test anomaly \\
\midrule
Candle      & 900 & 100 & 100 \\
Capsules    & 542 & 60  & 100 \\
Cashew      & 450 & 50  & 100 \\
Chewinggum  & 453 & 50  & 100 \\
Fryum       & 450 & 50  & 100 \\
Pipe fryum  & 450 & 50  & 100 \\
\midrule
Total       & 3245 & 360 & 600 \\
\bottomrule
\end{tabular}
\end{table}

All methods evaluate 1, 2, 4, and 8 normal references with seeds 0, 1, and 2.
The same reference filenames are used for the template baseline, official
AnomalyDINO, and multi-layer improvement. The first 20 normal images per
category are excluded from these references.

\subsection{Implementation Details}
The backbone is the pretrained DINOv2 ViT-S/14 checkpoint with a 384-dimensional
patch embedding and a smaller-edge size of 448 pixels. The backbone is frozen.
Nearest-neighbor search uses exact L2 distance on normalized features and runs
on CPU FAISS; the encoder runs on one NVIDIA RTX A6000. We use a batch size of
one, matching the official implementation, and report whether feature
extraction is included in timing.

The prediction outputs are evaluated per category and then averaged with equal
category weight. Image AUROC is the main metric. Pixel AUROC and AUPRO use all
test images, including normal images with zero masks. Because full-resolution
TIFF storage exceeded the safe free-space budget on the shared server, official
AnomalyDINO maps are evaluated at their native $32\times32$ patch grid and masks
are downsampled with nearest-neighbor interpolation. This is the prediction
resolution actually produced by the model and is identical across the official
and improved variants.

\subsection{Efficiency Protocol}
Peak GPU memory is measured after resetting CUDA statistics and loading the
model. Mean inference time is read from per-image measurements and includes
image loading, resizing, patch encoding, PCA masking, nearest-neighbor search,
and score aggregation, but excludes one-time model loading. We also benchmark
the encoder over 50 synchronized forward passes on a fixed $448\times448$
image.
